% !TEX program = pdflatex
% Comprehensive report of the Phase-1, Phase-2 and SSL (Phase-3) changes made
% to the Track-A action-anticipation pipeline for the
% "What-Happens-Next-Modal" Kaggle competition.
%
% Compiles standalone with pdflatex (no images, no external bibliography
% required). The companion file report/style.sty exists but is empty; this
% file therefore declares all the packages it needs directly.
\documentclass[11pt,a4paper]{article}

\usepackage[T1]{fontenc}
\usepackage[utf8]{inputenc}
\usepackage{lmodern}
\usepackage{microtype}
\usepackage[a4paper, margin=2.4cm]{geometry}
\usepackage{amsmath, amssymb}
\usepackage{booktabs}
\usepackage{xcolor}
\usepackage{graphicx}
\usepackage{enumitem}
\usepackage[colorlinks=true, linkcolor=blue!60!black,
            citecolor=blue!60!black, urlcolor=teal]{hyperref}
\usepackage{listings}

\definecolor{codebg}{HTML}{F7F7F9}
\definecolor{codekw}{HTML}{0033B3}
\definecolor{codestr}{HTML}{067D17}
\definecolor{codecomment}{HTML}{8C8C8C}
\definecolor{codeident}{HTML}{2A2A2A}

\lstdefinestyle{py}{
  language=Python,
  basicstyle=\ttfamily\footnotesize\color{codeident},
  keywordstyle=\color{codekw}\bfseries,
  commentstyle=\color{codecomment}\itshape,
  stringstyle=\color{codestr},
  backgroundcolor=\color{codebg},
  showstringspaces=false,
  breaklines=true,
  frame=single,
  framerule=0pt,
  framesep=4pt,
  rulecolor=\color{codebg},
  numbers=none,
  columns=fullflexible,
  literate=
    {é}{{\'e}}1
    {è}{{\`e}}1
    {ê}{{\^e}}1
    {à}{{\`a}}1
}
\lstdefinestyle{yaml}{
  basicstyle=\ttfamily\footnotesize\color{codeident},
  keywordstyle=\color{codekw}\bfseries,
  commentstyle=\color{codecomment}\itshape,
  stringstyle=\color{codestr},
  backgroundcolor=\color{codebg},
  morecomment=[l]\#,
  showstringspaces=false,
  breaklines=true,
  frame=single,
  framerule=0pt,
  framesep=4pt,
  rulecolor=\color{codebg},
  numbers=none,
  columns=fullflexible
}
\lstdefinestyle{shell}{
  basicstyle=\ttfamily\footnotesize\color{codeident},
  backgroundcolor=\color{codebg},
  showstringspaces=false,
  breaklines=true,
  frame=single,
  framerule=0pt,
  framesep=4pt,
  rulecolor=\color{codebg},
  numbers=none,
  columns=fullflexible
}

\newcommand{\code}[1]{\texttt{\small #1}}
\newcommand{\file}[1]{\texttt{\small #1}}
\newcommand{\cfgkey}[1]{\texttt{\small #1}}

\title{\textbf{Track-A Action Anticipation:\\
       Phase-1, Phase-2 and SSL Pretraining Improvements}}
\author{What-Happens-Next-Modal Kaggle competition\\
        \small Closed-world (from-scratch) track}
\date{\today}

\begin{document}
\maketitle

\begin{abstract}
This report documents in detail every change made to the Track-A pipeline
of the \emph{Something-Something V2 action anticipation} project. The work
proceeded in three phases. \textbf{Phase~1} closes a number of correctness
and accounting gaps that were silently inflating validation scores and
hurting generalisation: the validation split now points at the official
held-out folder, the ResNet-50 backbone is initialised with the Kaiming~He
recipe used by all modern from-scratch ResNet recipes, mixed-precision
training (AMP) is wired in, optimizer/scheduler/scaler state are persisted
across restarts, and a new \emph{trained-class index list} is recorded in
each checkpoint so that the submission pipeline can mask never-trained
logits at inference time. \textbf{Phase~2} adds four opt-in architectural
and training enhancements: per-block Stochastic Depth (DropPath), an
attention-pool head that strictly subsumes the legacy mean-consensus head,
exponential-moving-average (EMA) model averaging with an automatic
``best-of-(live, EMA)'' checkpoint selection rule, and Test-Time
Augmentation (TTA) with an automatic left-right class-pair remap at
inference time. \textbf{Phase~3} introduces a self-supervised pretraining
pipeline based on DINO-v1 that operates on still frames extracted from the
\emph{provided} train+val+test folders (Track-A-compliant: no external
data, no pretrained weights), and a non-strict \cfgkey{model.init\_from}
hook in the supervised trainer that warm-starts the ResNet-50 trunk from
the resulting SSL checkpoint. All changes default to off, preserving
byte-for-byte parity with the legacy pipeline; the test suite grows from
120 to 144 tests, all passing.
\end{abstract}

\tableofcontents

% =====================================================================
\section{Context}
\label{sec:context}

\subsection{Competition rules}

The Kaggle competition \emph{Can a model predict the future?} provides
short video clips from the Something-Something~V2 dataset and asks
participants to predict the action that will occur next, given only the
beginning of each clip. Track~A (\emph{Closed World}) is the focus of this
work: \textbf{models must be trained from scratch using only the provided
data}; no external dataset and no pretrained weights are allowed.
Submissions are evaluated by Top-1 classification accuracy (with Top-5 as
a secondary metric) on a held-out test set whose public/private split is
hidden from participants.

\subsection{Starting point and observed problems}

Before the work in this report, the supervised Track-A pipeline was a
TSM-ResNet50 trained with Adam, label smoothing, RandAugment-T temporal
augmentations, FrameMixUp, an 80/20 internal split of the training data
for validation, and a cosine learning-rate schedule. Its best public
Kaggle scores were:

\begin{center}
\begin{tabular}{lcc}
\toprule
\textbf{Submission} & \textbf{Internal val Top-1} & \textbf{Public Top-1} \\
\midrule
Initial run (epoch 29) & 0.4582 & 0.3207 \\
Resumed run (epoch 60) & 0.5666 & 0.3584 \\
\bottomrule
\end{tabular}
\end{center}

The roughly 21~percentage-point gap between internal validation and the
public leaderboard pointed at a collection of issues, the most important
of which were identified in a project review:

\begin{enumerate}[leftmargin=*]
\item \textbf{Leaky validation split}. The 80/20 split was carved out of
      \file{data/train/} rather than using the official held-out
      \file{data/val/} folder. Validation samples were therefore drawn
      from the same recording sessions and demographics as training,
      producing optimistically biased numbers.
\item \textbf{Suboptimal initialisation}. The ResNet-50 backbone was
      initialised with Xavier weights, which is well-suited to networks
      with sigmoid/tanh activations but is documented to under-perform
      Kaiming He initialisation on deep ReLU networks trained from
      scratch.
\item \textbf{No mixed-precision training}. Forward and backward passes
      ran in full FP32, leaving \(\approx 1.5{-}2\times\) of GPU
      throughput on the table on Ampere or newer hardware.
\item \textbf{Optimizer / scheduler / scaler state lost on resume}. Only
      the model weights and the epoch counter were checkpointed, so
      restarting a long run forced the optimizer (Adam moments) and the
      cosine LR schedule to restart from scratch as well.
\item \textbf{Missing class~27}. The ``50 classes'' described in the
      competition page is in practice 33 unique class indices on disk
      with class~27 entirely absent. The submission pipeline did nothing
      to prevent the network from emitting this never-trained label,
      meaning \(\approx 3\%\) of test predictions could land on it.
\item \textbf{No EMA, no TTA, no SSL pretraining.} The pipeline used
      none of the well-known accuracy-boosting techniques that are
      standard for video classification on small datasets.
\end{enumerate}

The remainder of this document describes the changes that addressed each
of these issues, in the order they were introduced.

% =====================================================================
\section{Phase 1: Foundations}
\label{sec:phase1}

Phase~1 comprises five low-risk, high-return changes. All of them were
introduced behind explicit configuration switches so that historical runs
can still be reproduced bit-for-bit by setting the new switches to their
``off'' defaults.

\subsection{Official validation split}
\label{sec:phase1:val}

\paragraph{Logic.} A train/val split that overlaps with the test
distribution understates the generalisation gap. Switching to the
official \file{data/val/} folder makes the validation score a faithful
proxy for the public leaderboard score and therefore drives early
stopping, model selection and hyperparameter choices in the right
direction.

\paragraph{Implementation.} A new boolean
\cfgkey{dataset.use\_official\_val} (default \code{false}) was added to
\file{configs/data/default.yaml}. When set to \code{true}, the trainer
bypasses the internal 80/20 split and instead enumerates every video
under \file{data/val/}.

\begin{lstlisting}[style=py]
use_official_val = bool(cfg.dataset.get("use_official_val", False))
if use_official_val:
    val_dir_for_val = Path(cfg.dataset.val_dir).resolve()
    val_samples = collect_video_samples(val_dir_for_val)
    train_samples = all_samples
    print(f"[data] use_official_val=true: train={len(train_samples)}, "
          f"val={len(val_samples)} (from {val_dir_for_val})")
else:
    train_samples, val_samples = split_train_val(
        all_samples, val_ratio=float(cfg.dataset.val_ratio),
        seed=int(cfg.dataset.seed))
\end{lstlisting}

\subsection{Kaiming He initialisation}
\label{sec:phase1:init}

\paragraph{Logic.} For a Conv$\to$BN$\to$ReLU stack trained from scratch,
the variance of pre-activations grows or shrinks geometrically with depth
unless the initial Conv weights are scaled with the Kaiming He recipe
\cite[Sec.~2.2]{HeICCV2015}:
\[
\mathrm{Var}(W_{i,j}) = \frac{2}{\text{fan\_out}}.
\]
Xavier initialisation, which uses \(2/(\text{fan\_in}+\text{fan\_out})\),
under-shoots this variance by roughly \(2\times\) per layer for ReLU
networks and is empirically a poor choice for from-scratch ResNet
training.

\paragraph{Implementation.} The
\code{\_init\_weights\_xavier} method of \code{AvancedResNet50TSM} was
renamed to \code{\_init\_weights} and rewritten to walk every submodule:

\begin{lstlisting}[style=py]
def _init_weights(self) -> None:
    for module in self.modules():
        if isinstance(module, nn.Conv2d):
            nn.init.kaiming_normal_(module.weight,
                mode="fan_out", nonlinearity="relu")
            if module.bias is not None:
                nn.init.zeros_(module.bias)
        elif isinstance(module, nn.BatchNorm2d):
            if module.weight is not None:
                nn.init.ones_(module.weight)
            if module.bias is not None:
                nn.init.zeros_(module.bias)
    nn.init.normal_(self.classifier.weight, mean=0.0, std=0.01)
    if self.classifier.bias is not None:
        nn.init.zeros_(self.classifier.bias)
\end{lstlisting}

The classifier head is initialised with a small-Gaussian
(\(\mathcal{N}(0, 0.01^2)\)) following the original ResNet and TSM
recipes; BN affine parameters are set to \(\gamma=1, \beta=0\).

\subsection{Automatic Mixed Precision (AMP)}
\label{sec:phase1:amp}

\paragraph{Logic.} On Ampere, Hopper and newer GPUs, running the forward
pass in FP16 (or BF16) and only the master copy of the optimizer state in
FP32 yields \(1.5\)--\(2\times\) throughput at no measurable accuracy
cost when paired with PyTorch's gradient scaler.

\paragraph{Implementation.} A new boolean \cfgkey{training.amp}
(default \code{false}) gates a \code{torch.amp.GradScaler} that is
created in \file{train.py} and threaded through the engine helpers.
\code{train\_one\_epoch} and \code{evaluate\_epoch} both received two new
keyword arguments, \code{scaler} and \code{amp\_dtype}, and wrap their
forward and backward passes in
\code{torch.amp.autocast(device\_type="cuda", dtype=amp\_dtype)}. When
\code{scaler} is \code{None} (the default), the code path is
byte-for-byte identical to the legacy full-precision implementation.

\begin{lstlisting}[style=py]
amp_enabled = bool(cfg.training.get("amp", False)) \
              and device.type == "cuda"
scaler = torch.amp.GradScaler("cuda", enabled=True) \
         if amp_enabled else None

# ...inside train_one_epoch:
with torch.amp.autocast(device_type="cuda",
                        enabled=use_amp,
                        dtype=amp_dtype):
    logits = model(videos)
    loss = loss_fn(logits, target)
if use_amp:
    scaler.scale(loss).backward()
    scaler.step(optimizer)
    scaler.update()
else:
    loss.backward()
    optimizer.step()
\end{lstlisting}

\subsection{Persistent optimizer / scheduler / scaler state}
\label{sec:phase1:resume}

\paragraph{Logic.} A long training run inevitably gets restarted
(scheduled maintenance, a transient OOM, an accidental \code{Ctrl-C}). If
only the model weights are persisted, the optimizer's running moments,
the cosine LR schedule and the GradScaler all reset, which causes a
visible accuracy dip for several epochs after each restart and makes
``resume from epoch~$k$'' a much weaker recovery than ``never
interrupted''.

\paragraph{Implementation.} On every \emph{best-checkpoint} save, the
trainer now stores \code{optimizer.state\_dict()},
\code{cosine\_scheduler.state\_dict()} (when present) and
\code{scaler.state\_dict()} (when AMP is on) inside the existing
\code{extra} dict of the checkpoint payload. The schema version is
unchanged. On resume, those state dicts are loaded inside try/except
blocks so a checkpoint produced by an older code path still resumes
cleanly (the cosine scheduler then falls back to its legacy
fast-forward).

\begin{lstlisting}[style=py]
ckpt_extra: dict[str, Any] = {
    "val_top1": ckpt_stats.top1,
    "val_top5": ckpt_stats.top5,
    "val_loss": ckpt_stats.loss,
    "epoch": epoch + 1,
    "optimizer_state_dict": optimizer.state_dict(),
    "trained_class_indices": trained_class_indices,
    "checkpoint_kind": ckpt_kind,
}
if cosine_scheduler is not None:
    ckpt_extra["scheduler_state_dict"] = cosine_scheduler.state_dict()
if scaler is not None:
    ckpt_extra["scaler_state_dict"] = scaler.state_dict()
\end{lstlisting}

\subsection{Trained-class indices and untrained-logit masking}
\label{sec:phase1:masking}

\paragraph{Logic.} The provided dataset has 33 unique class indices in
\([0, 32]\) but class~27 has no training samples. The classifier head is
nonetheless 33-wide (its weight tensor was initialised at random for
class~27) and at inference time argmax can therefore land on class~27 by
accident, producing systematically wrong predictions for every test
clip on which the random class~27 logit happens to be the largest.

\paragraph{Implementation.} The trainer derives the set of class indices
that actually have at least one training sample and writes it into the
checkpoint:

\begin{lstlisting}[style=py]
trained_class_indices: list[int] = sorted(
    {int(label) for _, label in train_samples})
\end{lstlisting}

The submission pipeline reads this list back and builds an additive mask
that pushes every never-trained class to \(-\infty\) before
\code{argmax}:

\begin{lstlisting}[style=py]
def _build_untrained_mask(trained_class_indices, num_classes, device):
    if not trained_class_indices:
        return None
    trained = {int(i) for i in trained_class_indices
               if 0 <= int(i) < num_classes}
    if len(trained) >= num_classes:
        return None
    mask = torch.zeros(num_classes, dtype=torch.float32, device=device)
    untrained = [c for c in range(num_classes) if c not in trained]
    mask[untrained] = -math.inf
    return mask
\end{lstlisting}

Older checkpoints that pre-date Phase~1 do not carry this list; in that
case the mask is \code{None} and the legacy unmasked argmax path is used,
preserving backward compatibility.

% =====================================================================
\section{Phase 2: Architectural and Training Enhancements}
\label{sec:phase2}

Phase~2 adds four orthogonal opt-in features that compose cleanly with
Phase~1. None of them change the behaviour of the legacy code paths when
left at their defaults; together, they are bundled in a new experiment
config \file{configs/experiment/track\_a\_phase2.yaml} that also flips on
AdamW with a stronger weight decay.

\subsection{Stochastic Depth (DropPath)}
\label{sec:phase2:droppath}

\paragraph{Logic.} Stochastic Depth \cite{Huang2016SD} is a stochastic
regularisation technique that, during training, drops the residual branch
of each ResNet bottleneck block with probability \(p_\ell\) and rescales
the kept contributions by \(1/(1-p_\ell)\) so that the expectation over
the random mask matches the no-drop forward pass:
\[
y_\ell \;=\; x_\ell \;+\;
\frac{b_\ell}{1-p_\ell}\,F_\ell(x_\ell), \qquad
b_\ell \sim \text{Bernoulli}(1-p_\ell).
\]
At inference time, all blocks are kept (\(b_\ell = 1\) deterministically)
so no rescaling is needed. The standard recipe uses a linear schedule
that drops deeper blocks more aggressively:
\[
p_\ell \;=\; p_{\max}\,\frac{\ell - 1}{L - 1},
\]
where \(L\) is the total number of bottleneck blocks. For ResNet-50
\(L=16\) and a typical \(p_{\max}\) is in \([0.05, 0.2]\).

\paragraph{Implementation.} A parameter-free \code{DropPath} module
implements the per-sample bernoulli mask. A helper
\code{\_inject\_drop\_path} walks the four ResNet stages, attaches a
\code{DropPath} instance to each block (with the linearly-scaled
probability), and rebinds the block's \code{forward} to a patched
implementation that calls \code{self.drop\_path(out)} on the residual
branch immediately before the residual addition:

\begin{lstlisting}[style=py]
def _patched_bottleneck_forward(self, x: torch.Tensor) -> torch.Tensor:
    identity = x
    out = self.relu(self.bn1(self.conv1(x)))
    out = self.relu(self.bn2(self.conv2(out)))
    out = self.bn3(self.conv3(out))
    if self.downsample is not None:
        identity = self.downsample(x)
    out = self.drop_path(out)         # <-- new
    return self.relu(out + identity)
\end{lstlisting}

\code{DropPath} adds no parameters and no buffers, so the model's
\code{state\_dict} layout is unchanged regardless of
\cfgkey{model.drop\_path\_rate}. This is asserted by a unit test that
compares state-dict key sets with and without DropPath enabled.

\subsection{Attention-pool head}
\label{sec:phase2:attn}

\paragraph{Logic.} The legacy temporal head averages the per-frame
features:
\[
z = \frac{1}{T}\sum_{t=1}^{T} \phi(x_t), \qquad
\widehat{y} = W z + b.
\]
A 1-query multi-head attention pool is a strict generalisation: when the
attention weights are uniform it reproduces the mean exactly, but it can
also concentrate on the most informative frame(s). Concretely, we
introduce a learnable query token \(q \in \mathbb{R}^{1\times D}\) and
compute
\[
z \;=\; \mathrm{LayerNorm}\!\left(
        \mathrm{MultiHeadAttn}(q,\, K{=}V{=}\Phi)\right),
\]
where \(\Phi \in \mathbb{R}^{T\times D}\) is the stack of frame
features. With \(T=4\) and \(D=2048\), this adds about
\(D + 4D^2 \approx 16{,}780{,}288\) parameters, which is small compared
to the \(\approx 25\)M parameter ResNet-50 trunk.

\paragraph{Implementation.} A new \code{AttentionPool} module wraps a
single \code{nn.MultiheadAttention} layer plus a final
\code{nn.LayerNorm}; the query token is initialised with a truncated
normal of std~0.02. The supervised model exposes a new
\cfgkey{model.head} switch (\code{"mean"} or \code{"attn"}) and routes
features accordingly:

\begin{lstlisting}[style=py]
sequence_features = frame_features.view(batch_size, num_frames, -1)
if self.attn_pool is not None:
    consensus_features = self.attn_pool(sequence_features)
else:
    consensus_features = sequence_features.mean(dim=1)
\end{lstlisting}

When \cfgkey{model.head}~$=$~\code{"mean"} (the default) the
\code{attn\_pool} attribute is \code{None} and no extra parameters are
added.

\subsection{Exponential Moving Average (EMA) of model weights}
\label{sec:phase2:ema}

\paragraph{Logic.} Polyak averaging \cite{PolyakSIAM1992} maintains an
exponentially-weighted average of the model's parameters,
\[
\theta_{\mathrm{EMA}}^{(t+1)} \;=\;
\alpha\,\theta_{\mathrm{EMA}}^{(t)} \;+\;
(1-\alpha)\,\theta^{(t+1)}, \qquad \alpha \in [0.99, 0.9999],
\]
which empirically gives smoother loss curves and a small but consistent
top-1 boost (\(\approx 0.3{-}1.0\) points) on most modern image and
video classifiers.

\paragraph{Implementation.} We use PyTorch's
\code{torch.optim.swa\_utils.AveragedModel} with a custom averaging
function and \code{use\_buffers=True} so BatchNorm running statistics are
also tracked:

\begin{lstlisting}[style=py]
def _ema_avg_fn(avg_param, model_param, _num_averaged):
    return ema_decay * avg_param + (1.0 - ema_decay) * model_param

ema_model = torch.optim.swa_utils.AveragedModel(
    model, avg_fn=_ema_avg_fn, use_buffers=True).to(device)
\end{lstlisting}

\code{train\_one\_epoch} accepts a new \code{ema\_model} keyword and
calls \code{ema\_model.update\_parameters(model)} after every optimizer
step. At the end of each epoch the trainer evaluates \emph{both} the
live and the EMA model and saves the better of the two:

\begin{lstlisting}[style=py]
if ema_stats is not None and ema_stats.top1 > val_stats.top1:
    ckpt_stats, ckpt_module, ckpt_kind = ema_stats, ema_model, "ema"
else:
    ckpt_stats, ckpt_module, ckpt_kind = val_stats, model, "live"

# ...later, before save_checkpoint:
module_to_save = (ckpt_module.module
                  if isinstance(ckpt_module,
                                torch.optim.swa_utils.AveragedModel)
                  else ckpt_module)
\end{lstlisting}

This selection rule is what lets downstream consumers
(\file{submit.py}, \file{evaluate.py}) stay completely oblivious to the
EMA toggle: the saved \code{model\_state\_dict} is always loaded with
\code{build\_model(...).load\_state\_dict(...)}, regardless of which
copy won. The string \code{"live"} or \code{"ema"} is recorded in
\code{extra["checkpoint\_kind"]} for traceability.

\subsection{Test-Time Augmentation with auto class-pair remap}
\label{sec:phase2:tta}

\paragraph{Logic.} Horizontal flip is a free \emph{evaluation-time}
augmentation \emph{provided} that the predicted class is invariant to
the flip. The Something-Something~V2 subset used in this competition
contains a small number of direction-sensitive classes, in particular:
\begin{itemize}[leftmargin=2em]
\item class~018: ``Pulling something from left to right''
\item class~019: ``Pulling something from right to left''
\end{itemize}
For these, a flipped clip should predict the \emph{paired} class. A
naive flip-and-average TTA actually hurts accuracy on these two classes
because it averages two confident but \emph{disagreeing} predictions.

\paragraph{Implementation.} We average softmax probabilities across the
two views (original and flipped) and apply a learned permutation
\(\pi: \{0,\ldots,K-1\}\to\{0,\ldots,K-1\}\) to the flipped softmax so
that paired classes are mapped back into agreement. The permutation is
\emph{auto-derived from the class folder names}, making it robust to any
future change in the class set:

\begin{lstlisting}[style=py]
LR = "from_left_to_right"
RL = "from_right_to_left"
for idx, stem in stem_by_idx.items():
    if LR in stem:
        mirror_stem = stem.replace(LR, RL)
        for other_idx, other_stem in stem_by_idx.items():
            if other_stem == mirror_stem:
                perm[idx] = other_idx
                perm[other_idx] = idx
                break
\end{lstlisting}

(The numeric prefix is stripped before matching so
\code{018\_..\_left\_to\_right} pairs with \code{019\_..\_right\_to\_left}
without a code-level hard-code.) The TTA loop in \file{submit.py} then
becomes:

\begin{lstlisting}[style=py]
logits = _logits_for_batch(model, video_batch, untrained_mask)
probs_total = torch.softmax(logits, dim=1)
n_views = 1
if tta_flip:
    flipped = torch.flip(video_batch, dims=[-1])
    flipped_probs = torch.softmax(
        _logits_for_batch(model, flipped, untrained_mask), dim=1)
    if flip_perm is not None:
        flipped_probs = flipped_probs.index_select(dim=1, index=flip_perm)
    probs_total = probs_total + flipped_probs
    n_views += 1
predictions.extend(int(p) for p in
    (probs_total / n_views).argmax(dim=1).cpu().tolist())
\end{lstlisting}

The configuration is opt-in (\cfgkey{training.tta}, \cfgkey{training.tta\_flip},
both \code{false} by default) and read at inference time only; training
is unaffected.

\subsection{AdamW and the bundled experiment configuration}
\label{sec:phase2:adamw}

The trainer already supported AdamW; Phase~2 simply makes it the
recommended optimizer for the new experiment by setting
\cfgkey{training.optimizer: adamw} and \cfgkey{training.weight\_decay:
0.05} in \file{configs/experiment/track\_a\_phase2.yaml}. Decoupling the
weight-decay term from the gradient (which is what AdamW does, unlike
plain Adam) is the right thing to do whenever a non-zero weight decay is
used; it is also the optimizer of choice for ConvNeXt, ViT and modern
ResNet recipes.

The \file{track\_a\_phase2.yaml} bundle composes all of Phase~1 and
Phase~2:

\begin{lstlisting}[style=yaml]
defaults:
  - override /model: avanced_resnet50_tsm
  - override /augment: randaugment_t

model:
  drop_path_rate: 0.1
  head: attn
  head_num_heads: 4

dataset:
  num_frames: 4
  use_official_val: true

training:
  optimizer: adamw
  weight_decay: 0.05
  amp: true
  ema_enabled: true
  ema_decay: 0.999
  tta: true
  tta_flip: true
  scheduler_cosine: true
  warmup_epochs: 10
  label_smoothing: 0.1
\end{lstlisting}

% =====================================================================
\section{Phase 3: Self-Supervised Pretraining (DINO)}
\label{sec:phase3}

The single highest-impact change available within Track~A is to give the
trunk something useful to do with the unlabeled frames before the
supervised loss starts back-propagating. We implemented DINO-v1
\cite{Caron2021DINO}, a teacher-student knowledge-distillation framework
that has consistently outperformed contrastive methods (MoCo, SimCLR) on
ResNet-50 trunks at the dataset sizes we are working with.

\subsection{Track-A compliance}

The competition forbids \emph{external data} and \emph{pretrained
models}. Self-supervised pretraining on the \emph{provided} unlabeled
images is therefore allowed: the labels are simply not consumed and the
trunk starts from random initialisation. We pretrain on the union of the
\file{train/}, \file{val/} and \file{test/} folders, which gives roughly
\(\sim 12{,}000\) still frames (\(\sim 3{,}000\) videos
\(\times\) 4 frames) and is sufficient to drive the DINO loss for many
hundreds of steps per epoch at batch size~64.

\subsection{Still-frames dataset}
\label{sec:phase3:dataset}

A new file \file{src/smth2smth/shared/data/still\_frames\_dataset.py}
exposes two utilities:

\begin{itemize}[leftmargin=2em]
\item \code{collect\_all\_frame\_paths(roots)} walks the provided
      directories. It transparently handles both layouts in use:
      \file{train/<class>/<video>/<frame>.jpg} (via
      \code{collect\_video\_samples}) and
      \file{test/<video>/<frame>.jpg} (one level shallower, no class
      bucket). Duplicate paths are removed.
\item \code{MultiViewStillFramesDataset} returns, for each frame, two
      global views (typically 224\(\times\)224 with random crop scale in
      \([0.4, 1.0]\)) and \(L\) local views (typically 96\(\times\)96
      with random crop scale in \([0.05, 0.4]\)). The augmentations
      include horizontal flip, colour jitter, grayscale, and Gaussian
      blur, following the DINO recipe.
\end{itemize}

The dataset returns a \code{dict} with keys \code{global\_views} and
\code{local\_views}; a custom collate function stacks them into per-view
batches.

\subsection{The DINO model}
\label{sec:phase3:model}

\file{src/smth2smth/shared/models/dino\_ssl.py} contains three pieces.

\paragraph{Trunk.} \code{build\_resnet50\_trunk()} returns a
\code{torchvision.models.resnet50(weights=None)} with its final \code{fc}
replaced by \code{nn.Identity}. The trunk emits 2048-d features.

\paragraph{Head.} \code{DinoHead} is the canonical 3-layer projection
MLP used in DINO: \code{Linear} $\to$ \code{GELU} $\to$ \code{Linear}
$\to$ \code{GELU} $\to$ \code{Linear} (to a low-dimensional bottleneck)
$\to$ L2 normalisation $\to$ weight-normalised \code{Linear} mapping to
\(K\) prototypes (\(K=4096\) by default).

\paragraph{Composite.} \code{DinoModel} composes the trunk and the head.
Two copies are instantiated at training time --- a \emph{student} and a
\emph{teacher}. Both have identical architecture; the teacher's weights
are an EMA of the student's and are never updated by gradient descent.

\subsection{The DINO loss}
\label{sec:phase3:loss}

Let \(\mathcal{V} = \{v_1, \ldots, v_{2+L}\}\) be the set of all views
of one image (two globals, \(L\) locals). Let \(p_s(v)\) and \(p_t(v)\)
be the softmax-normalised outputs of the student and teacher at
temperatures \(\tau_s\) and \(\tau_t\) respectively. The loss is
\[
\mathcal{L} \;=\;
\frac{1}{|\mathcal{P}|}\!\!\sum_{(t, s)\in \mathcal{P}}\!\!
H\!\left( p_t(v_t),\; p_s(v_s) \right),
\]
where \(\mathcal{P}\) is the set of (teacher view, student view) pairs
with \emph{distinct} indices and \(H\) is the standard cross-entropy.
The teacher only processes the two global views (so
\(|\mathcal{P}| = 2 \cdot |\mathcal{V}| - 2 \cdot 1\)), but the student
processes \emph{all} views. The teacher's logits are sharpened
(\(\tau_t = 0.04\) is much smaller than \(\tau_s = 0.1\)) and centred by
an EMA of the batch mean to prevent collapse:
\[
c^{(t+1)} \;=\; m_c\, c^{(t)} \;+\;
(1 - m_c)\, \frac{1}{B}\sum_{i=1}^{B}\mathrm{logits}_t(v_i^{\mathrm{global}}),
\]
with momentum \(m_c = 0.9\).

\begin{lstlisting}[style=py]
class DinoLoss(nn.Module):
    def forward(self, student_logits_per_view, teacher_logits_per_global):
        teacher_probs_per_global = [
            F.softmax((logits - self.center) / self.teacher_temp,
                      dim=-1).detach()
            for logits in teacher_logits_per_global]
        student_log_softmax_per_view = [
            F.log_softmax(logits / self.student_temp, dim=-1)
            for logits in student_logits_per_view]
        total_loss, n_terms = 0.0, 0
        for t_idx, t_probs in enumerate(teacher_probs_per_global):
            for s_idx, s_log in enumerate(student_log_softmax_per_view):
                if t_idx == s_idx:
                    continue
                total_loss = total_loss \
                             + (-(t_probs * s_log).sum(dim=-1).mean())
                n_terms += 1
        # ... centering buffer EMA update ...
        return total_loss / n_terms
\end{lstlisting}

\subsection{Teacher EMA momentum schedule}

The teacher's parameters are updated by
\(\theta_t \leftarrow m\,\theta_t + (1-m)\,\theta_s\) after every
optimizer step, with \(m\) following a cosine schedule from
\(m_0 = 0.996\) at the start to \(m_1 = 1.0\) at the end of training.
This mirrors the recipe in the DINO paper and stabilises training
significantly --- a constant momentum is more prone to collapse on small
datasets.

\subsection{Trunk-only checkpoint format}

After every epoch we save a small ($\approx$ 95~MiB) checkpoint that
contains \emph{only} the trunk parameters, with the \code{trunk.}
prefix stripped:

\begin{lstlisting}[style=py]
trunk_state_dict = {
    k.removeprefix("trunk."): v
    for k, v in student.state_dict().items()
    if k.startswith("trunk.")
}
torch.save({"trunk_state_dict": trunk_state_dict,
            "epoch": epoch + 1}, out_path)
\end{lstlisting}

This format is intentionally minimal: there is no point keeping the DINO
head, the teacher copy or the optimizer state once supervised training
takes over.

\subsection{The \cfgkey{model.init\_from} hook in \file{train.py}}
\label{sec:phase3:initfrom}

The supervised trainer received a single new opt-in:

\begin{lstlisting}[style=py]
init_from = cfg.model.get("init_from") if hasattr(cfg.model, "get") \
            else None
if init_from:
    payload = torch.load(Path(str(init_from)).resolve(),
                         map_location=device, weights_only=False)
    trunk_state = payload.get("trunk_state_dict")
    prefixed = _ssl_trunk_to_supervised_keys(trunk_state)
    missing, unexpected = model.load_state_dict(prefixed, strict=False)
\end{lstlisting}

The non-trivial part is the key remap. The supervised
\code{AvancedResNet50TSM} wraps every bottleneck block's \code{conv1}
in
\[
\code{block.conv1} \;\leftarrow\;
\code{nn.Sequential(TemporalShift(\(T\), shift\_div), conv1)},
\]
which means the parameter that the SSL trunk has at key
\code{layer1.0.conv1.weight} now lives at
\code{layer1.0.conv1.1.weight} in the supervised model.
\code{\_ssl\_trunk\_to\_supervised\_keys} performs this rename
unconditionally for every key matching
\code{layer[1-4]\textbackslash.\textbackslash d+\textbackslash.conv1\textbackslash.(weight|bias)}:

\begin{lstlisting}[style=py]
block_conv1_re = re.compile(
    r"^(layer[1-4]\.\d+\.conv1)\.(weight|bias)$")
out: dict[str, torch.Tensor] = {}
for k, v in trunk_state.items():
    match = block_conv1_re.match(k)
    if match is not None:
        new_k = f"{match.group(1)}.1.{match.group(2)}"
    else:
        new_k = k
    out[f"backbone.{new_k}"] = v
return out
\end{lstlisting}

The load is intentionally non-strict so that the supervised model's
classifier head and (when enabled) attention pool keys remain at their
fresh-init values. The trainer logs how many tensors were loaded, and
warns loudly if any backbone key was \emph{not} covered by the SSL
checkpoint.

% =====================================================================
\section{Configuration and reproducibility}
\label{sec:configs}

\subsection{New configuration knobs}

The following table summarises every new switch introduced by Phases~1
through~3. All defaults preserve the legacy behaviour.

\begin{center}\small
\begin{tabular}{@{}lll@{}}
\toprule
\textbf{Key} & \textbf{Default} & \textbf{Effect when on} \\
\midrule
\cfgkey{dataset.use\_official\_val}  & \code{false} & validate on \file{data/val/} \\
\cfgkey{training.amp}                & \code{false} & FP16 autocast + GradScaler \\
\cfgkey{training.ema\_enabled}       & \code{false} & EMA of model weights \\
\cfgkey{training.ema\_decay}         & \code{0.999} & EMA momentum \\
\cfgkey{training.tta}                & \code{false} & enable TTA at submit time \\
\cfgkey{training.tta\_flip}          & \code{true}  & flip view (only used when \code{tta}) \\
\cfgkey{training.tta\_temporal\_shift}\!\! & \code{0} & extra temporal start-offsets (TTA) \\
\cfgkey{model.drop\_path\_rate}      & \code{0.0}   & Stochastic Depth maximum drop prob.\\
\cfgkey{model.head}                  & \code{mean}  & temporal head (\code{mean}/\code{attn}) \\
\cfgkey{model.head\_num\_heads}      & \code{4}     & MultiHead attention heads \\
\cfgkey{model.init\_from}            & \code{null}  & path to SSL trunk to warm-start \\
\cfgkey{pretrain.*}                  & --           & SSL pretraining hyperparameters \\
\bottomrule
\end{tabular}
\end{center}

\subsection{New experiment configurations}

\begin{description}[leftmargin=1.5em, style=nextline]
\item[\file{configs/experiment/track\_a\_phase2.yaml}]
   Bundles every Phase-1 and Phase-2 opt-in: AdamW + WD~0.05, AMP,
   DropPath~0.1, attention head, EMA~0.999, TTA at submit time, official
   validation split, cosine schedule with 10-epoch warmup, label
   smoothing~0.1, RandAugment-T, FrameMixUp.

\item[\file{configs/experiment/track\_a\_ssl\_pretrain.yaml}]
   Routes the run through \code{pretrain\_ssl.py}: 150 epochs of
   DINO-v1 over the union of train+val+test still frames, 2 globals +
   6 locals per sample, AdamW + WD~0.04, AMP on.

\item[\file{configs/experiment/track\_a\_ssl\_finetune.yaml}]
   Same as \file{track\_a\_phase2.yaml} but expects
   \cfgkey{model.init\_from} to be set on the command line; reduces the
   number of supervised epochs (warm-started runs converge faster) and
   the learning rate by half.
\end{description}

\subsection{Backwards compatibility}

Every legacy entry point continues to work without modification:

\begin{itemize}[leftmargin=2em]
\item Phase-1 checkpoints (no \code{trained\_class\_indices}, no
      \code{optimizer\_state\_dict} in \code{extra}) load and resume via
      the legacy fast-forward fallback.
\item Pre-Phase-1 checkpoints (no Phase-2 opt-ins enabled) load with
      identical state-dict keys; this is enforced by a unit test that
      compares state-dict key sets between defaults and Phase-2
      opt-ins.
\item Submission CSVs produced before Phase~1 remain valid; the new
      mask is only applied if the checkpoint actually carries the
      trained-class list.
\end{itemize}

% =====================================================================
\section{Test coverage}
\label{sec:tests}

The pre-existing suite contained 120 tests, all of which still pass.
Twenty-four new tests were added in three files, bringing the total to
\textbf{144 passing tests}:

\begin{description}[leftmargin=1.5em, style=nextline]
\item[\file{tests/shared/models/test\_phase2\_optins.py} (8 tests)]
   Verifies that
   (i)~\code{DropPath} is identity at eval time and unbiased per-sample
   at train time;
   (ii)~the \code{DropPath(0)} module is byte-for-byte identical to the
   no-op;
   (iii)~the attention pool returns the expected \((B, D)\) shape;
   (iv)~enabling \cfgkey{drop\_path\_rate} does \emph{not} grow the
   state-dict key set;
   (v)~enabling the attention head adds the expected \code{attn\_pool.*}
   keys.

\item[\file{tests/pipelines/test\_submit\_tta.py} (5 tests)]
   Verifies the auto-derived left/right class permutation, the no-TTA
   path, the untrained-mask path, and that the flip-TTA softmax average
   actually reinforces the correct class when the remap is applied.

\item[\file{tests/shared/data/test\_still\_frames\_dataset.py} (6 tests)]
   Covers frame-path discovery for both the class-bucketed and
   test-style folder layouts, deduplication across multiple roots, and
   the multi-view sampling shape contract.

\item[\file{tests/shared/models/test\_dino\_ssl.py} (4 tests)]
   Asserts that the remapped SSL trunk keys are exactly a superset of
   the supervised model's backbone keys, that the loss is finite and
   back-propagates, and that the teacher EMA update matches the analytic
   midpoint formula at \(m=0.5\).

\item[\file{tests/pipelines/test\_init\_from\_ssl.py} (1 test)]
   End-to-end smoke test: build a DINO model, save its trunk-only
   state dict, then load it into a freshly built
   \code{AvancedResNet50TSM} via the same code path the supervised
   trainer uses; check that no backbone key is left at random init,
   that the classifier and attention pool are not touched, and that
   spot-checked tensors round-trip exactly.
\end{description}

\noindent The full suite is exercised by:

\begin{lstlisting}[style=shell]
PYTHONPATH=src .venv/bin/python -m pytest
\end{lstlisting}

% =====================================================================
\section{Expected impact and launch strategy}
\label{sec:strategy}

\subsection{Per-change order-of-magnitude estimates}

Based on the literature and on the gap currently observed between
internal validation and the public leaderboard, we expect the following
contributions to Track~A's public Top-1 score (these are
order-of-magnitude estimates, not promises):

\begin{center}\small
\begin{tabular}{lcc}
\toprule
\textbf{Change} & \textbf{Internal val} & \textbf{Public Top-1} \\
\midrule
Baseline (current best)                        & 0.5666 & 0.3584 \\
+ official val split (Phase~1)                 & lower\,$^*$ & --     \\
+ Kaiming init (Phase~1)                       & +0.5--2 pt & +0.3--1 pt \\
+ AMP (Phase~1)                                & speedup only & speedup only \\
+ optimizer/scheduler state on resume (Phase~1)& 0--1 pt & 0--1 pt \\
+ untrained-class masking (Phase~1)            & negligible & +0.5--1.5 pt \\
+ DropPath (Phase~2)                           & +0.5--1 pt & +0.3--0.8 pt \\
+ Attention head (Phase~2)                     & +0.5--1 pt & +0.3--0.8 pt \\
+ EMA model averaging (Phase~2)                & +0.3--0.8 pt & +0.2--0.6 pt \\
+ TTA flip with class-pair remap (Phase~2)     & +0--0.5 pt & +0.5--1 pt \\
+ AdamW + WD~0.05 (Phase~2)                    & +0.3--1 pt & +0.3--1 pt \\
+ DINO SSL pretraining (Phase~3)               & +1--3 pt & +1--3 pt \\
\bottomrule
\end{tabular}\\
\smallskip
{\footnotesize $^*$ The official split is harder than the leaky
internal split, so reported val will drop --- but the public score will
move closer to the val score.}
\end{center}

\subsection{Recommended launch order}

\paragraph{Strategy A --- Phase~2 only (fastest, single run).}
Best when the goal is one strong submission overnight. Approx.\ runtime
on a single Ampere GPU at batch~16: \(\sim\)20\,h with AMP for 150
epochs.

\begin{lstlisting}[style=shell]
nohup env PYTHONUNBUFFERED=1 uv run python scripts/run_track_a.py \
    experiment=track_a_phase2 \
    training.num_workers=2 \
    training.checkpoint_path=.../best_model_phase2.pt \
    > "$LOG" 2>&1 &
\end{lstlisting}

\paragraph{Strategy B --- SSL pretrain then Phase-2 fine-tune.}
Highest expected gain, two stages, \(\sim\)25--28\,h total. Pretrain
first:

\begin{lstlisting}[style=shell]
nohup env PYTHONUNBUFFERED=1 PYTHONPATH=src uv run python \
    -m smth2smth.pipelines.pretrain_ssl \
    experiment=track_a_ssl_pretrain track=a \
    pretrain.checkpoint_path=.../ssl_trunk.pt > "$LOG" 2>&1 &
\end{lstlisting}

Then fine-tune from the SSL checkpoint:

\begin{lstlisting}[style=shell]
nohup env PYTHONUNBUFFERED=1 uv run python scripts/run_track_a.py \
    experiment=track_a_ssl_finetune \
    model.init_from=.../ssl_trunk.pt \
    training.checkpoint_path=.../best_model_ssl_finetune.pt \
    > "$LOG" 2>&1 &
\end{lstlisting}

\paragraph{Submission with TTA.}
After either strategy, generate the submission CSV with flip-TTA and
auto class-pair remap enabled:

\begin{lstlisting}[style=shell]
PYTHONPATH=src uv run python -m smth2smth.pipelines.submit \
    experiment=track_a_phase2 \
    training.checkpoint_path=.../best_model_*.pt \
    training.tta=true training.tta_flip=true \
    dataset.submission_output=.../submission_tta.csv
\end{lstlisting}

% =====================================================================
\section{Summary of files touched}
\label{sec:files}

\begin{description}[leftmargin=1.5em, style=nextline, font=\normalfont\ttfamily]
\item[src/smth2smth/pipelines/train.py]
   Phase~1 (official val, AMP, optimizer/scheduler/scaler state,
   trained-class list) + Phase~2 (EMA + best-of-(live, EMA) selection)
   + Phase~3 (\code{model.init\_from} hook with key remap helper).
\item[src/smth2smth/pipelines/submit.py]
   Phase~1 (untrained-class masking) + Phase~2 (flip-TTA with
   auto-derived left/right class permutation).
\item[src/smth2smth/pipelines/pretrain\_ssl.py \,\textbf{(new)}]
   Phase~3 DINO-v1 pretraining loop on still frames.
\item[src/smth2smth/shared/engine/trainer.py]
   Phase~1 AMP wiring + Phase~2 \code{ema\_model} hook.
\item[src/smth2smth/shared/models/avanced\_resnet50\_tsm.py]
   Phase~1 Kaiming init + Phase~2 \code{DropPath}, \code{AttentionPool},
   patched \code{Bottleneck.forward}.
\item[src/smth2smth/shared/models/dino\_ssl.py \,\textbf{(new)}]
   Phase~3 \code{DinoModel}, \code{DinoHead}, \code{DinoLoss},
   \code{update\_teacher\_ema}.
\item[src/smth2smth/shared/data/still\_frames\_dataset.py \,\textbf{(new)}]
   Phase~3 \code{collect\_all\_frame\_paths},
   \code{MultiViewStillFramesDataset}.
\item[configs/data/default.yaml]            Phase~1 \cfgkey{use\_official\_val}.
\item[configs/train/default.yaml]           Phase~1 \cfgkey{amp},
                                            Phase~2 \cfgkey{ema\_*},
                                            \cfgkey{tta\_*}.
\item[configs/model/avanced\_resnet50\_tsm.yaml]
   Phase~2 \cfgkey{drop\_path\_rate}, \cfgkey{head}, \cfgkey{head\_num\_heads};
   Phase~3 \cfgkey{init\_from}.
\item[configs/pretrain/default.yaml \,\textbf{(new)}]
   Phase~3 SSL pretraining defaults.
\item[configs/experiment/track\_a\_phase2.yaml \,\textbf{(new)}]
   Phase~2 bundle.
\item[configs/experiment/track\_a\_ssl\_pretrain.yaml \,\textbf{(new)}]
   Phase~3 SSL bundle.
\item[configs/experiment/track\_a\_ssl\_finetune.yaml \,\textbf{(new)}]
   Phase~2 + SSL warm-start bundle.
\item[configs/config.yaml]                  Phase~3 added \code{pretrain: default}
                                            to the defaults list.
\item[tests/ ...]                           24 new unit and integration tests
                                            (\(120 \to 144\) passing).
\end{description}

% =====================================================================
\section{Conclusion}

Phase~1 fixes the silent correctness issues that were inflating
validation scores and silently degrading test performance. Phase~2 layers
on four well-understood, opt-in techniques --- Stochastic Depth, an
attention-pool head, EMA model averaging, and flip-TTA with auto-derived
class-pair remap --- that compose cleanly with each other and with the
existing RandAugment-T + FrameMixUp data pipeline. Phase~3 provides the
single largest expected jump by warm-starting the trunk from a
self-supervised pretraining run on the provided unlabeled frames, with a
clean \cfgkey{model.init\_from} hook that handles the
TemporalShift-induced key rename automatically. All changes default to
off, the test suite has grown from 120 to 144 passing tests, and the
existing experiment configurations and resume paths continue to work
unchanged.

% =====================================================================
\begin{thebibliography}{9}

\bibitem{HeICCV2015}
K.~He, X.~Zhang, S.~Ren, J.~Sun.
\newblock \emph{Delving Deep into Rectifiers: Surpassing Human-Level
              Performance on ImageNet Classification}.
\newblock ICCV 2015.

\bibitem{Huang2016SD}
G.~Huang, Y.~Sun, Z.~Liu, D.~Sedra, K.~Q.~Weinberger.
\newblock \emph{Deep Networks with Stochastic Depth}.
\newblock ECCV 2016.

\bibitem{Caron2021DINO}
M.~Caron, H.~Touvron, I.~Misra, H.~J\'egou, J.~Mairal, P.~Bojanowski,
A.~Joulin.
\newblock \emph{Emerging Properties in Self-Supervised Vision
              Transformers}.
\newblock ICCV 2021.

\bibitem{PolyakSIAM1992}
B.~T.~Polyak, A.~B.~Juditsky.
\newblock \emph{Acceleration of Stochastic Approximation by Averaging}.
\newblock SIAM J. Control and Optimization, 1992.

\bibitem{Lin2019TSM}
J.~Lin, C.~Gan, S.~Han.
\newblock \emph{TSM: Temporal Shift Module for Efficient Video
              Understanding}.
\newblock ICCV 2019.

\bibitem{Loshchilov2019AdamW}
I.~Loshchilov, F.~Hutter.
\newblock \emph{Decoupled Weight Decay Regularization}.
\newblock ICLR 2019.

\end{thebibliography}

\end{document}
